\begin{statementbox}
对于椭圆 $C: \frac{x^2}{a^2} + \frac{y^2}{b^2} = 1$ ($a>b>0$), 设 $l$ 是椭圆的任意一条切线. 关于距离关系有以下核心结论:
\begin{enumerate}
    \item \textbf{原心距公式}: 原点 $O(0, 0)$ 到切线 $y = kx \pm \sqrt{a^2k^2 + b^2}$ 的距离 $d = \frac{\sqrt{a^2k^2 + b^2}}{\sqrt{k^2 + 1}}$;
    \item \textbf{焦点距离积定值}: 若左右焦点分别为 $F_1, F_2$, 则 $F_1$ 到切线 $l$ 的距离 $d_1$ 与 $F_2$ 到切线 $l$ 的距离 $d_2$ 之积满足 $d_1 \cdot d_2 = b^2$;
    \item \textbf{垂足位置性质}: 焦点在切线上的投影 (垂足) 必落在椭圆的主圆 $x^2 + y^2 = a^2$ 上.
\end{enumerate}
\end{statementbox}